\clearpage
\chapter{Result and Discussion}
\label{chap:results}

\section{Performance Evaluation Overview}

The evaluation of Savior Glass covers six functional areas: Taka detection, Taka authentication, emotion recognition, text recognition, object detection and processing speed. Each module was evaluated on held-out data that was not used for training or model selection, using the protocol described in Section~\ref{sec:method}. Model training and evaluation were run on laptop GPUs (NVIDIA RTX~4060 for the Taka detector, RTX~3050 for the other models); speed on the laptop CPU is also reported. Every number in this chapter was copied from a saved evaluation file of the project. When a measurement does not exist, it is written as \textit{not measured} rather than estimated.

\section{Taka Detection Performance}

\subsection{Training Behaviour}
The YOLOv8s detector was fine-tuned from COCO weights for 80 epochs at 640$\times$640 pixels with batch size 16, SGD (initial learning rate 0.01, momentum 0.937, weight decay 0.0005, three warm-up epochs), mosaic 1.0 and mixup 0.15. Horizontal and vertical flips were disabled because a mirrored note is not a real note. Training converged quickly: mAP@0.5 was above 0.97 after 8 epochs and above 0.99 after 11 epochs, while mAP@0.5:0.95 kept improving slowly as the boxes became tighter (Figure~\ref{fig:traincurves}).

\begin{figure}[H]
    \centering
    \includegraphics[width=\textwidth]{fig_training_curves.png}
    \caption{Training curves of the YOLOv8s Taka detector: box, class and DFL losses (top) and mAP@0.5, mAP@0.5:0.95, precision and recall (bottom) against epoch}
    \label{fig:traincurves}
\end{figure}

\subsection{Results on Held-Out Composites}
Table~\ref{tab:det_overall} reports the final performance on the 2,238 held-out test composites.

\begin{table}[!htbp]
\centering
\caption{Taka detection performance on the synthetic test set (2,238 images)}
\label{tab:det_overall}
\begin{tabular}{|l|c|}
\hline
\textbf{Metric} & \textbf{Value} \\ \hline
Precision & 0.998 \\ \hline
Recall & 0.997 \\ \hline
mAP@0.5 & 0.995 \\ \hline
mAP@0.5:0.95 & 0.847 \\ \hline
Inference time (RTX 4060 laptop GPU) & 7.3 ms / image ($\approx$135 FPS) \\ \hline
Pre-process / post-process time & 1.7 ms / 1.7 ms \\ \hline
\end{tabular}
\end{table}

A later re-validation of the same weights on the locally available validation data gave mAP@0.5 = 0.994, mAP@0.5:0.95 = 0.844, precision 0.991 and recall 0.992, which agrees with the original numbers. The gap between mAP@0.5 (0.995) and mAP@0.5:0.95 (0.847) shows that the detector almost always finds and correctly names the note, but that the box edges are less precise at strict IoU thresholds. One plausible reason is the slight anti-aliasing at the pasted, warped note borders, which does not correspond to a sharp physical edge.

\begin{table}[!htbp]
\centering
\caption{Per-denomination detection metrics on the synthetic test set}
\label{tab:det_perclass}
\begin{tabular}{|l|c|c|c|c|c|}
\hline
\textbf{Class} & \textbf{Instances} & \textbf{Precision} & \textbf{Recall} & \textbf{mAP@0.5} & \textbf{mAP@0.5:0.95} \\ \hline
2 Taka & 180 & 0.998 & 1.000 & 0.995 & 0.838 \\ \hline
5 Taka & 264 & 0.998 & 0.996 & 0.995 & 0.844 \\ \hline
10 Taka & 172 & 0.995 & 0.988 & 0.995 & 0.848 \\ \hline
20 Taka & 192 & 0.998 & 1.000 & 0.995 & 0.847 \\ \hline
50 Taka & 172 & 1.000 & 0.993 & 0.995 & 0.851 \\ \hline
100 Taka & 228 & 0.999 & 1.000 & 0.995 & 0.845 \\ \hline
200 Taka & 172 & 0.999 & 1.000 & 0.995 & 0.843 \\ \hline
500 Taka & 500 & 1.000 & 1.000 & 0.995 & 0.853 \\ \hline
1000 Taka & 172 & 0.994 & 1.000 & 0.995 & 0.858 \\ \hline
\textbf{All} & \textbf{2,052} & \textbf{0.998} & \textbf{0.997} & \textbf{0.995} & \textbf{0.847} \\ \hline
\end{tabular}
\end{table}

Results are uniformly strong across the nine classes (Table~\ref{tab:det_perclass}), including the smallest source class (10 Taka, 419 source images), so the model does not simply favour the most frequent class (500 Taka, 1,243 source images). Because mAP@0.5 rounds to 0.995 for every class, mAP@0.5:0.95 (0.838--0.858) is the more informative column for comparing denominations. The normalised confusion matrix (Figure~\ref{fig:det_confusion}) is almost perfectly diagonal; the remaining errors are missed or spurious detections against the background, not confusions between two denominations. For a currency reader this is the safer failure mode: the glass may ask the user to hold the note again, but it rarely names the wrong value.

\begin{figure}[H]
    \centering
    \includegraphics[width=0.7\textwidth]{fig_confusion_matrix.png}
    \caption{Normalised confusion matrix of the Taka detector on the synthetic test set}
    \label{fig:det_confusion}
\end{figure}

\begin{figure}[H]
    \centering
    \includegraphics[width=0.95\textwidth]{fig_qualitative.jpg}
    \caption{Example predictions on held-out composites with predicted class and confidence. The note is localised across different scales, rotations and backgrounds.}
    \label{fig:qualitative}
\end{figure}

\subsection{Results on Raw Phone Photographs (Source Domain, Not Held Out)}
Table~\ref{tab:wild} shows the result of the same detector on the 450 raw phone photographs (50 per denomination). A note counts as \textit{detected} if any box is predicted, and as \textit{correct} if the top-scoring box has the true denomination.

\begin{table}[!htbp]
\centering
\caption{Taka detection on 450 raw phone photographs (50 per denomination). Source-domain check: these photographs come from the detector's training source and are not held out. Independent results: Table~\ref{tab:external}.}
\label{tab:wild}
\begin{tabular}{|l|c|c|c|p{4.6cm}|}
\hline
\textbf{True class} & \textbf{Detected} & \textbf{Correct} & \textbf{Top-1 (\%)} & \textbf{Errors} \\ \hline
2 Taka & 49 & 47 & 94.0 & 2 read as 5 Taka, 1 missed \\ \hline
5 Taka & 50 & 50 & 100.0 & -- \\ \hline
10 Taka & 48 & 48 & 96.0 & 2 missed \\ \hline
20 Taka & 50 & 50 & 100.0 & -- \\ \hline
50 Taka & 48 & 46 & 92.0 & 1 as 5 Taka, 1 as 500 Taka, 2 missed \\ \hline
100 Taka & 48 & 48 & 96.0 & 2 missed \\ \hline
200 Taka & 50 & 50 & 100.0 & -- \\ \hline
500 Taka & 50 & 50 & 100.0 & -- \\ \hline
1000 Taka & 50 & 47 & 94.0 & 3 read as 500 Taka \\ \hline
\textbf{Overall} & \textbf{443 (98.4\%)} & \textbf{436} & \textbf{96.9} & \textbf{7 missed, 7 wrong value} \\ \hline
\end{tabular}
\end{table}

On raw photographs the detector finds 98.4\% of the notes and names 96.9\% correctly. Because these photographs come from the same source as the training foregrounds, this number does not measure transfer to new notes; the independent collections in Table~
ef{tab:external} do (91.5\% on Bangla Money). The errors are informative: 2 and 5 Taka are both small, similar-coloured notes, and the 1000 and 500 Taka notes share a similar layout and portrait. As explained in Chapter~\ref{chap:dataset}, some physical notes in this set may also occur in the training foregrounds, so this is an optimistic estimate. The 450 images were drawn from the raw folder without excluding the photographs used as training foregrounds, so this table is a source-domain check, not a test; the independent results below replace it as the real-photograph result.

\subsection{Results on Independent Photographs}
Two public collections that were never used for training or selection were scored with the same detector at confidence 0.25 (Table~\ref{tab:external}). A third, ``A Diverse Image Dataset for Bangladeshi Currency Recognition'', was excluded because its files are byte-identical to the training source.

\begin{table}[!htbp]
\centering
\caption{Taka detection on independent photograph collections}
\label{tab:external}
\begin{tabular}{|l|c|p{7cm}|}
\hline
\textbf{Set} & \textbf{n} & \textbf{Result} \\ \hline
Bangla Money (Kaggle), 8 denominations & 1,536 & 91.5\% correct denomination; no note found 3.9\%; 95.2\% correct when a note is found \\ \hline
NSTU-BDTAKA detection test & 186 & class-agnostic AP@0.5 0.796; precision 0.832, recall 0.694 \\ \hline
NSTU-BDTAKA hand-held close-ups & 1,144 & 18.5\% correct denomination; no note found 66.9\% \\ \hline
Bangla Money 1-Taka notes (unknown class) & 101 & 57 not detected; 44 announced as a known note (37 as 5 Taka) \\ \hline
\end{tabular}
\end{table}

The detector transfers to full-note photographs taken by other people, with 1000 Taka the weakest class (78.4\%, mostly read as 500 Taka). It does not transfer to close-ups in which fingers cover part of a folded note, and it has no ``unknown note'' output.

\subsection{Fall-Back Classifier}
The MobileNetV3-Small classifier used as a fall-back reached a best validation accuracy of 99.2\% on the nine-class BanglaTaka crops. This is a whole-note classification score on clean crops, not a detection score in a scene.

\subsection{Comparison with Related Work}
Compared with published detectors (Table~\ref{tab:currency_compare}), our mAP@0.5:0.95 of 0.847 on composites is higher than the 0.804 reported for a custom-head YOLOv8n multi-currency model \cite{yolov8nse}, and our precision and recall (0.998/0.997) are higher than the 0.95/0.93 of the YOLOv5 baseline trained on 3,111 manually annotated images \cite{nstubdtaka2024}. Our labelled dataset is about seven times larger and required no manual annotation. However, the test sets are different and our headline scores are on synthetic images, so the comparison indicates competitiveness, not superiority on a common benchmark.

\section{Taka Authentication Performance}

\subsection{Accuracy for One to Six Views}
All authentication models were evaluated on the same 208 test notes with the first $k$ views, for $k = 1,\dots,6$. Table~\ref{tab:auth_views} compares the CNN+ViT baseline, which was trained with all six views, with the prefix-robust model (PRMVT).

\begin{table}[!htbp]
\centering
\caption{Authentication accuracy on 208 test notes with the first $k$ views (seed 42)}
\label{tab:auth_views}
\begin{tabular}{|c|c|c|c|}
\hline
\textbf{Views} & \textbf{CNN+ViT baseline} & \textbf{PRMVT (proposed)} & \textbf{Gain (points)} \\ \hline
1 & 73.56\% & \textbf{97.12\%} & +23.6 \\ \hline
2 & 87.02\% & \textbf{98.08\%} & +11.1 \\ \hline
3 & 91.35\% & \textbf{98.08\%} & +6.7 \\ \hline
4 & 91.83\% & \textbf{98.08\%} & +6.3 \\ \hline
5 & 89.90\% & \textbf{98.56\%} & +8.7 \\ \hline
6 & 91.83\% & \textbf{98.08\%} & +6.3 \\ \hline
\end{tabular}
\end{table}

The baseline is much weaker with one view (73.6\%) than with six (91.8\%). An earlier version of our fusion model trained only on six views showed the same problem even more strongly: 96.6\% with six views but only 68.8\% with one view. This is the view-count distribution shift described in Section~\ref{sec:method}. The prefix-robust protocol removes it: PRMVT stays between 97.1\% and 98.6\% for every number of views. The difference to the baseline is statistically significant under McNemar's test at one view ($p = 4.3\times10^{-11}$) and at six views ($p = 0.0020$). After a Bonferroni correction over 36 algorithm-by-view comparisons, the one-, two-, three- and five-view gains remain significant, while the four- and six-view gains do not; this correction is reported so that the six-view improvement is not overstated. The PRMVT numbers in this chapter were re-evaluated after a numerical fault was fixed: the entropy of a fully confident prediction was computed as $0 \cdot \log 0$, which is undefined, and the model then reported a probability of 0.5 for that note. The fault affected 23 of the 208 test notes at six views and made the earlier figures slightly lower (six-view accuracy 97.6\% instead of 98.1\%).

\begin{figure}[H]
    \centering
    \includegraphics[width=0.85\textwidth]{fig_auth_views.png}
    \caption{Authentication accuracy against the number of views for the CNN+ViT baseline, PRMVT and the other training variants tried in this project}
    \label{fig:auth_views}
\end{figure}

\subsection{Repeatability over Seeds}
PRMVT was trained with three random seeds (42, 43, 44). The mean accuracy was 96.5\% (standard deviation 1.6 points) with one view and 98.1\% (standard deviation 1.0 point) with six views; the lowest mean over the six view counts was 96.5\% (Table~\ref{tab:multiseed}). The result is therefore not a lucky seed, although three seeds is a small sample.

\begin{table}[!htbp]
\centering
\caption{PRMVT accuracy over three seeds (mean $\pm$ standard deviation)}
\label{tab:multiseed}
\begin{tabular}{|c|c|c|c|c|c|c|}
\hline
\textbf{Views} & 1 & 2 & 3 & 4 & 5 & 6 \\ \hline
\textbf{Mean (\%)} & 96.47 & 98.40 & 98.08 & 97.60 & 98.56 & 98.08 \\ \hline
\textbf{SD (points)} & 1.55 & 0.56 & 0.48 & 0.83 & 0.48 & 0.96 \\ \hline
\end{tabular}
\end{table}

\subsection{Other Metrics and Calibration}
For the six-view operating point of the fixed-count fusion model, the improvement over the baseline is consistent across metrics: accuracy 0.966 vs.\ 0.918, macro-F1 0.966 vs.\ 0.914, balanced accuracy 0.965 vs.\ 0.906, ROC-AUC 0.989 vs.\ 0.971, and expected calibration error (ECE) 0.030 vs.\ 0.079. A lower ECE means that the confidence reported by the model matches its real accuracy better, which is important if the glass tells the user ``probably counterfeit''. For PRMVT, the six-view ECE without any calibration is 0.015, so its confidence is already well matched to its accuracy; calibration maps fitted on the validation set (temperature scaling and a quality-aware map) do not improve on it (0.017--0.018). An earlier uncalibrated value of 0.061 was caused by the numerical fault described above, which reported a probability of 0.5 for notes the model was certain about. The confusion matrices (Figure~\ref{fig:auth_confusion}) show that PRMVT mainly removes the baseline's errors on counterfeit notes that were accepted as genuine (a fraction of 0.07 of all test notes for the baseline, 0.01 for PRMVT).

\begin{figure}[H]
    \centering
    \includegraphics[width=0.75\textwidth]{fig_auth_confusion.png}
    \caption{Normalised confusion matrices of the baseline and PRMVT at six views (0 = counterfeit, 1 = genuine)}
    \label{fig:auth_confusion}
\end{figure}

\subsection{Robustness to Image Corruptions}
The models were tested under 13 image corruptions at several strengths (Figure~\ref{fig:auth_robust}). PRMVT is almost unaffected by Gaussian and motion blur, JPEG compression, scale change and sensor noise. It is clearly sensitive to low light, strong brightness, rotation, perspective and especially occlusion: at an occlusion of 55\% of the image, six-view accuracy fell to 51.9\%, which is worse than the baseline (81.3\%). Fine-tuning with low-light augmentation raised accuracy at the darkest setting (brightness factor 0.2) to 93.75\% while keeping 98.6\% on clean images, but a first, short occlusion fine-tune did not fix occlusion (68.3\%, with clean accuracy falling to 90.9\%). An analysis on the validation set showed why: under occlusion PRMVT called most genuine notes counterfeit, because a covered region looks like a missing security feature, and the earlier fine-tunes had kept the CNN frozen and trained only briefly. A longer fine-tune in which every view is randomly covered and the last CNN blocks can adapt, with the epoch chosen on the validation set, raised accuracy at 55\% occlusion to 88.5\% (87.0\% $\pm$ 2.2 over three seeds, 88.9\% for their average) while clean accuracy stayed at 98--99\%. This is still below 90\%. If the model is allowed to answer only when it is at least 99\% confident and otherwise asks the user to reposition the note, it answers about half of the occluded notes with 99.0\% accuracy, and wrong verdicts fall from 11.1\% to 0.5\% of occluded notes. These tests use a black rectangle; a real finger is neither black nor rectangular, so occlusion by fingers remains the main open failure mode.

\begin{figure}[H]
    \centering
    \includegraphics[width=\textwidth]{fig_auth_robustness.png}
    \caption{Accuracy drop (clean minus corrupted accuracy) at six views under 12 corruptions for PRMVT, the NDAL variant and the baseline}
    \label{fig:auth_robust}
\end{figure}

\subsection{Authentication on Whole-Note Photographs}
\label{sec:jaal_deployed}
The glass applies the authenticator to one crop of the note found by the detector, not to the close-up views it was trained on. This deployed path was scored on photographs that were not used for training or selection (Table~\ref{tab:jaal_deployed}).

\begin{table}[!htbp]
\centering
\caption{Deployed genuine/counterfeit verdict on whole-note photographs (one crop, threshold 0.5)}
\label{tab:jaal_deployed}
\begin{tabular}{|l|c|c|}
\hline
\textbf{Set} & \textbf{Notes judged} & \textbf{Wrong verdict (95\% Wilson CI)} \\ \hline
Bangla Money, genuine & 1,464 & 47.7\% (45.2--50.3) \\ \hline
BanglaTaka raw, genuine & 441 & 59.2\% (54.5--63.7) \\ \hline
Counterfeit-currency set, genuine & 1,187 & 20.4\% (18.2--22.8) \\ \hline
Counterfeit-currency set, counterfeit & 75 & 40.0\% (29.7--51.3) \\ \hline
\end{tabular}
\end{table}

A genuine note called counterfeit, or a counterfeit called genuine, is the most harmful error for a blind user. The glass therefore no longer speaks the verdict: it announces the denomination followed by ``the counterfeit check was not done''. The high close-up accuracy in Table~\ref{tab:auth_views} does not carry over to whole-note crops.

\subsection{Other Training Variants and Negative Results}
Fifteen other training variants were tried on the same split (Figure~\ref{fig:auth_views}). Several of them reach 95--98\% at one view, for example a focal-style loss variant (NDAL) \cite{focalloss} with 98.1\% at one view, so the accuracy alone does not single out one algorithm; the common factor is that all of them see varying numbers of views during training. Some variants failed: a simulated federated version stayed at 74.0\% with one view and below the baseline from two views on, and one retrained variant became worse than its first version. We also tried learned policies that choose which view to capture next and when to stop (rules based on image quality, confidence, diversity and uncertainty, and a cost-aware rule). On PRMVT, averaged over three seeds, none of them beat simply taking the views in their fixed order (96.5\% with one view and 97.6--98.6\% with two to six): the learned rules picked a weaker first view and reached 89.3--90.2\% with one view, and the cost-aware rule with early stopping stopped after one view at 89.3\%. A prefix stopping rule on the fixed order reached 97.1\% with an average of 1.01 views, which is simply because PRMVT is already accurate with one view. Earlier, much lower figures for these policies (down to 42\%) were caused by evaluation faults that have since been fixed, not by the policies themselves. These negative results are reported so that the method is not presented as better than it is.

\subsection{Speed}
On the RTX~3050 laptop GPU, one six-view authentication takes a median of 42.6~ms (95th percentile 69.5~ms). The model file is 14.5~MB. The time on the Raspberry Pi~5 has not been measured.

\section{Emotion Recognition Performance}

Several models were trained and compared on the official RAF-DB test set (Table~\ref{tab:emotion_runs}). The EfficientNet-V2-S model with flip test-time augmentation gave the best accuracy, 86.5\%, and is used in the system.

\begin{table}[!htbp]
\centering
\caption{Emotion recognition models compared during development}
\label{tab:emotion_runs}
\begin{tabular}{|p{4.8cm}|p{4.0cm}|c|c|}
\hline
\textbf{Model} & \textbf{Test set} & \textbf{Accuracy} & \textbf{Used} \\ \hline
FER+ ONNX & FER-2013 hold-out, n = 5,741 & 57.9\% & Fall-back \\ \hline
MobileNetV3-Large & RAF-DB official test, n = 3,068 & 81.0\% & Replaced \\ \hline
ViT (face-expression) fine-tune & RAF-DB official test & 84.3\% (peak) & No \\ \hline
EfficientNet-B2 (AffectNet) fine-tune & RAF-DB official test & 85.2\% (then over-fit) & No \\ \hline
\textbf{EfficientNet-V2-S + flip TTA} & \textbf{RAF-DB official test} & \textbf{86.5\%} & \textbf{Yes} \\ \hline
\end{tabular}
\end{table}

\begin{table}[!htbp]
\centering
\caption{Per-class results of the EfficientNet-V2-S emotion model on RAF-DB (n = 3,068)}
\label{tab:emotion_perclass}
\begin{tabular}{|l|c|c|c|c|}
\hline
\textbf{Class} & \textbf{Support} & \textbf{Precision} & \textbf{Recall} & \textbf{F1} \\ \hline
Happy & 1,185 & 0.946 & 0.925 & 0.936 \\ \hline
Sad & 478 & 0.839 & 0.839 & 0.839 \\ \hline
Angry & 162 & 0.789 & 0.809 & 0.799 \\ \hline
Surprise & 329 & 0.867 & 0.872 & 0.870 \\ \hline
Fear & 74 & 0.774 & 0.554 & 0.646 \\ \hline
Disgust & 160 & 0.671 & 0.588 & 0.627 \\ \hline
Neutral & 680 & 0.815 & 0.890 & 0.851 \\ \hline
\textbf{Macro average} & \textbf{3,068} & \textbf{0.815} & \textbf{0.782} & \textbf{0.795} \\ \hline
\end{tabular}
\end{table}

The accuracy is 86.5\% and the macro-F1 is 0.795 (Table~\ref{tab:emotion_perclass}). Happy, surprise, neutral and sad are recognised well; fear and disgust are the weakest classes because they have the fewest training images and are often confused with sad and neutral (Table~\ref{tab:emotion_cm}). The published state of the art on this split is about 92\% (POSTER~V2 \cite{posterv2}), so our model is a reasonable but not state-of-the-art recogniser. For adaptive speech, the most important decision is whether a person looks stressed (sad, fear, anger, disgust) or not, which is less demanding than the full seven-class task.

\begin{table}[!htbp]
\centering
\caption{Confusion matrix of the EfficientNet-V2-S emotion model on the RAF-DB official test set (rows: true class, columns: predicted class)}
\label{tab:emotion_cm}
\begin{tabular}{|l|c|c|c|c|c|c|c|}
\hline
\textbf{True / Pred.} & \textbf{Hap.} & \textbf{Sad} & \textbf{Ang.} & \textbf{Sur.} & \textbf{Fear} & \textbf{Dis.} & \textbf{Neu.} \\ \hline
\textbf{Happy} & \textbf{1096} & 11 & 8 & 12 & 1 & 11 & 46 \\ \hline
\textbf{Sad} & 14 & \textbf{401} & 5 & 4 & 1 & 14 & 39 \\ \hline
\textbf{Angry} & 7 & 2 & \textbf{131} & 4 & 4 & 4 & 10 \\ \hline
\textbf{Surprise} & 6 & 8 & 3 & \textbf{287} & 5 & 5 & 15 \\ \hline
\textbf{Fear} & 3 & 8 & 5 & 11 & \textbf{41} & 2 & 4 \\ \hline
\textbf{Disgust} & 11 & 15 & 13 & 4 & 0 & \textbf{94} & 23 \\ \hline
\textbf{Neutral} & 21 & 33 & 1 & 9 & 1 & 10 & \textbf{605} \\ \hline
\end{tabular}
\end{table}

\section{Text Recognition Performance}

\subsection{OCR Error Rates}
OCR quality is measured by the character error rate (CER) and word error rate (WER), the edit distance between recognised text and ground truth divided by the length of the ground truth. Table~\ref{tab:ocr} shows the results of EasyOCR with our pre-processing on the 80-phrase test set.

\begin{table}[!htbp]
\centering
\caption{OCR error rates of EasyOCR (Bangla + English) with CLAHE pre-processing}
\label{tab:ocr}
\begin{tabular}{|l|c|c|c|}
\hline
\textbf{Language} & \textbf{Samples} & \textbf{CER} & \textbf{WER} \\ \hline
English & 40 & 3.2\% & 10.4\% \\ \hline
Bangla & 40 & 27.8\% & 57.9\% \\ \hline
\textbf{Overall} & \textbf{80} & \textbf{15.5\%} & \textbf{34.2\%} \\ \hline
\end{tabular}
\end{table}

English text is read almost perfectly (3.2\% CER), but Bangla text has a CER of 27.8\%: on average more than one character in four is wrong, and more than half of the words contain an error. Inspection of the outputs shows that most Bangla errors are in vowel signs (\textit{kar}) that are dropped or placed on the wrong consonant, and in conjunct consonants (\textit{juktakkhor}) that are split into separate letters. The consonant skeleton of a word is usually correct, so a listener can often still guess the word, but the result is not reliable enough for unfamiliar text. A lexicon-based correction was tried; because its word list included the test phrases, its near-zero error is not a valid measurement and is not reported as a result.

\subsection{Handwritten Bangla Character Recognition}
The CNN trained on Ekush reached 90.0\% test accuracy on 38,608 characters from 305 unseen writers (best validation accuracy 91.1\%). Accuracy was 94.3\% for digits, 94.2\% for modifiers, 93.1\% for vowels, 89.9\% for consonants and 87.7\% for conjuncts. The lower accuracy for conjuncts confirms that compound characters are the hardest part of Bangla recognition, which is also the main source of the OCR errors above.

\section{Object Detection Performance}
The COCO-pretrained YOLOv8s model reached mAP@0.5 = 0.604 on the 250-image COCO val2017 subset (77 classes present). Large, distinctive objects such as buses, laptops, suitcases, dogs and beds had AP@0.5 above 0.90, while small or rare objects were weaker. The object mode reduces the effect of false alarms by using a confidence threshold of 0.50 and by announcing at most three objects per scan. Precision of the spoken announcements in real indoor use has not yet been measured.

\section{Processing Speed}
Table~\ref{tab:speed} summarises the measured speed. On a laptop CPU (AMD Ryzen~7 5800H, all three formats measured in the same run), the Taka detector takes a median of 109~ms per 640$\times$640 image (9.2 FPS) with PyTorch weights, 139~ms with FP32 ONNX Runtime and 101~ms (9.9 FPS) with the INT8 ONNX model, which is also the smallest file. An earlier INT8 figure (199~ms, 11.5~MB) belonged to a defective export that detected no notes and has been replaced. None of these figures is a Raspberry Pi~5 measurement.

\begin{table}[!htbp]
\centering
\caption{Measured processing speed (Raspberry Pi~5 not measured)}
\label{tab:speed}
\begin{tabular}{|p{4.6cm}|p{4.2cm}|c|c|}
\hline
\textbf{Model} & \textbf{Hardware} & \textbf{Time} & \textbf{Size} \\ \hline
YOLOv8s Taka (PyTorch) & RTX 4060 laptop GPU & 7.3 ms & 22.5 MB \\ \hline
YOLOv8s Taka (PyTorch) & Laptop CPU (Ryzen 7 5800H) & 109 ms & 22.5 MB \\ \hline
YOLOv8s Taka (FP32 ONNX) & Laptop CPU (Ryzen 7 5800H) & 139 ms & 44.8 MB \\ \hline
YOLOv8s Taka (INT8 ONNX) & Laptop CPU (Ryzen 7 5800H) & 101 ms & 18.0 MB \\ \hline
PRMVT authenticator, 6 views & RTX 3050 laptop GPU & 42.6 ms (median) & 14.5 MB \\ \hline
All models & Raspberry Pi~5 & Not measured & -- \\ \hline
\end{tabular}
\end{table}

\section*{Calculation Methods and Full Working}

This section presents the formulas used in this chapter and a worked example for each main table.

\subsection*{1) Detection Precision, Recall and mAP (Table~\ref{tab:det_overall})}
With $TP$, $FP$ and $FN$ the numbers of true positives, false positives and false negatives at a given IoU threshold,
\[
\text{Precision} = \frac{TP}{TP+FP}, \qquad \text{Recall} = \frac{TP}{TP+FN}, \qquad
\text{IoU}(A,B) = \frac{|A\cap B|}{|A\cup B|}.
\]
The average precision (AP) of one class is the area under its precision--recall curve; mAP@0.5 is the mean AP over the 9 classes at IoU $\geq 0.5$, and mAP@0.5:0.95 is the mean over the ten IoU thresholds $0.50, 0.55, \dots, 0.95$.

\subsection*{2) Raw-Photograph Detection Rate and Top-1 Accuracy (Table~\ref{tab:wild})}
\[
\text{Detection rate} = \frac{\text{notes with at least one box}}{\text{all notes}} = \frac{450-7}{450} = \frac{443}{450} = 98.4\%,
\]
\[
\text{Top-1 accuracy} = \frac{\text{notes with correct top box}}{\text{all notes}} = \frac{436}{450} = 96.9\%,
\]
where $436 = 48+48+47+47+50+50+50+46+50$ is the sum of the correct notes over the nine denominations.

\subsection*{3) Authentication Accuracy and McNemar's Test (Table~\ref{tab:auth_views})}
With 208 test notes, one-view PRMVT accuracy $= 202/208 = 97.12\%$ and baseline accuracy $= 153/208 = 73.56\%$. McNemar's test uses only the discordant notes: $n_{01}$ (baseline correct, PRMVT wrong) and $n_{10}$ (baseline wrong, PRMVT correct):
\[
\chi^2 = \frac{(|n_{01}-n_{10}|-1)^2}{n_{01}+n_{10}},
\]
with one degree of freedom. For example, for the fixed-count fusion model at six views $n_{01}=3$ and $n_{10}=13$, giving $\chi^2 = 81/16 = 5.06$ and $p = 0.024$.

\subsection*{4) Expected Calibration Error}
Predictions are grouped into $M = 10$ confidence bins $B_m$:
\[
\text{ECE} = \sum_{m=1}^{M} \frac{|B_m|}{n}\,\bigl|\text{acc}(B_m) - \text{conf}(B_m)\bigr|.
\]

\subsection*{5) Emotion Macro-F1 (Table~\ref{tab:emotion_perclass})}
For each class $c$, $F1_c = 2P_cR_c/(P_c+R_c)$. For fear, $F1 = 2\times0.774\times0.554/(0.774+0.554) = 0.646$. Macro-F1 is the unweighted mean of the seven class F1 scores: $(0.936+0.839+0.799+0.870+0.646+0.627+0.851)/7 = 0.795$.

\subsection*{6) OCR Error Rates (Table~\ref{tab:ocr})}
\[
\text{CER} = \frac{S+D+I}{N},
\]
where $S$, $D$ and $I$ are the character substitutions, deletions and insertions needed to turn the recognised text into the ground truth, and $N$ is the number of ground-truth characters. WER is computed in the same way on words.


\section{Discussion}
The results answer the research questions of Section~\ref{sec:Research Question}.

\textbf{System design (RQ1).} A complete three-mode assistive glass was implemented with offline models, Bangla speech, physical buttons, automatic start at boot and graceful fall-backs. The design runs on a laptop and is packaged for the Raspberry Pi~5; its speed on the Pi still has to be measured.

\textbf{Synthetic data for currency (RQ2).} Training only on automatically labelled composites gave near-perfect scores on held-out composites and, more importantly, 91.5\% correct denominations on 1,536 independently collected photographs, though only 18.5\% on hand-held close-ups. The main lesson is that realistic backgrounds and physically motivated augmentation matter more than a special network architecture.

\textbf{Multi-view authentication (RQ3).} A model trained with a fixed number of views fails when fewer views are available; training on random view prefixes fixes this and keeps accuracy at about 97--98\% from one to six views on a note-disjoint test set, with low variation across seeds. For a blind user this means that one clear view of the note is usually enough, and more views give a small additional benefit. Occlusion remains the main weakness: an occlusion fine-tune reaches 88.5\% with 55\% of each view covered, below the 90\% target, and a confidence threshold can turn most of the remaining errors into requests to reposition the note.

\textbf{Bangla OCR (RQ4).} Off-the-shelf OCR reads English well but Bangla poorly (27.8\% CER). Vowel signs and conjunct characters cause most errors. This is the weakest module of the glass and the one that most needs improvement.

\textbf{Interaction design (RQ5).} Separating capture from reading, per-class cooldowns, queue clearing on mode change, mixed-script speech and vibration patterns were designed for users without sight. Whether they actually help can only be shown with a user study, which has not been done yet.

\section{Limitations}
The following limitations must be considered when interpreting the results:
\begin{enumerate}[noitemsep]
    \item \textbf{No Raspberry Pi~5 measurements.} Latency, frame rate, CPU and memory use, temperature and battery life on the Pi~5 have not been measured. All speed figures are from laptops.
    \item \textbf{No user study.} The system has not been tested with visually impaired participants; usability claims are design intentions, not measured outcomes.
    \item \textbf{Synthetic and possibly overlapping test data for currency.} Headline detection metrics are on composites, and the raw-photograph set shares its source with the training data. On independent photographs the detector scores 91.5\% on full notes but 18.5\% on hand-held close-ups.
    \item \textbf{No counterfeit verdict on the glass.} The authenticator is accurate on JaalTaka close-up views but wrong on 20--59\% of genuine whole-note photographs, so the glass does not announce a verdict.
    \item \textbf{Single capture collection for authentication.} JaalTaka has no camera or session identifiers, so generalisation to other phones, cameras and lighting is untested. The authenticator is a visible-light aid, not a substitute for bank equipment.
    \item \textbf{Weak Bangla OCR.} A CER of 27.8\% is too high for reliable reading of unfamiliar Bangla text.
    \item \textbf{Occlusion and low light.} Authentication degrades when the note is partly covered: 88.5\% at 55\% synthetic occlusion after occlusion fine-tuning (51.9\% before), and occlusion by real fingers has not been measured.
    \item \textbf{Generic objects.} The object mode uses the 80 COCO classes, which do not include many objects important for navigation (stairs, open drains, rickshaws) and gives no distance.
    \item \textbf{Emotion camera direction.} Emotion-adaptive speech was tested with a user-facing webcam; on the glass the camera faces outward.
    \item \textbf{Small evaluation sets.} The OCR set has 80 phrases and the object subset 250 images, so these numbers have wide uncertainty.
\end{enumerate}
